\section{Moyenne}
\setcounter{exemple}{0}

\begin{exemplebox}
    Cinq enfants ont apporté respectivement 3, 3, 5, 6 et 13 bonbons à la
    garderie. L'animatrice décide que chaque enfant devrait avoir le même nombre
    de bonbons. Elle met tous les bonbons en commun. \\

    Elle obtient \answer{30} bonbons. 
    (Calcul: \answer{$3+3+5+6+13=30$})\\

    Elle les redistribue équitablement entre les 5 enfants.

    Chacun obtient \answer{6} bonbons (Calcul: \answer{$30\div5=6$}), c'est la \bb{moyenne}.\\

    \highlight[RougeClair]{\bb{Interprétation}}: Si tous les enfants avaient
    chacun \answer{6} bonbons, le nombre total de bonbons resterait le même.

\end{exemplebox}

\begin{methodebox}
    Pour calculer la moyenne, on fait la somme de toutes les valeurs puis on
    divise par l'effectif total :
    \begin{center}
        $\text{moyenne} = \dfrac{\text{somme des valeurs}}{\text{effectif
        total}} $
    \end{center}
\end{methodebox}

\begin{exemplebox}
    Jean a 6 notes en maths : 18, 12, 15, 14, 20 et 17. Calcule sa moyenne

    \vspace{5cm}

    \highlight[RougeClair]{\bb{Interprétation}}: Si toutes les notes étaient
    égales à \answer{16}, le total des points resterait le même.
\end{exemplebox}


\begin{exercicebox}
    Sans calculatrice déterminer la moyenne de ces valeurs:
    \begin{itemize}
        \item 14; 6; 7
        \item 15; 20; 18; 22; 25
        \item 140; 170; 160; 130
    \end{itemize}
\end{exercicebox}
